%\subsection{Conclusion}
Read through Ihde’s thesis of non-neutrality, digital audio tools are never just utilitarian; the tools themselves mediate and materialize particular worldviews about sound, time, and the body. Technopoiesis reframes this inevitability as an opening: cultural epistemologies can be intentionally embedded into design such that tools co-construct other, alternative, sonic worlds \parencite{Guerra_Ostergaard_2017}. Black speculative traditions and sonic fictions provide the epistemic and imaginative reservoir for this form of modification and embedding. In doing so, offering diagrams for different temporalities, logics, and modes of audition. These can ultimately be transduced into interfaces, algorithms, and workflows. The lens of computational culture clarifies why this remaking is necessary: default parameters and data regimes are not neutral; they sediment histories of exclusion that require algorithmic justice in both method and artifact. In sum, the conceptual strands, the harmonic threads, of post-phenomenology, technopoiesis, Afrofuturism, sonic fiction, and computation converge on a single design imperative: remake the mediations through embodied modification.

Our methodological wager is that such remaking happens best where critical technical practice and research-creation meet. From Agre, we take the split stance and the discipline of reflective build, "keeping one foot in craft and one in critique," and treating impasses as cues for reframing from within the work \parencite[p.~155]{Agre_1998}. From Loveless, we inherit a commitment to form-as-world-making and a micropolitical pedagogy that composes livable worlds “bit by bit” in the university-as-site \parencite[p.~102]{Loveless_2019}. From Manning and Massumi, we take a speculative pragmatism: construct initial conditions that catalyze emergent concepts-in-action, and iterate technique as event \parencite[p.~89]{Manning_Massumi_2014}. Together, these commitments compose Sonic Speculocultural Technopoiesis as both lens and method: we will critically and creatively design and critique with practitioners, treating the design process itself as a site of inquiry where Black sonic practices are not only represented but enacted in code, interface, and performance \parencite{Guerra_Ostergaard_2017}. This is how we will proceed. In step, by folding reflection into making so that critique becomes sound, and sound itself becomes a working theory.